% year: תשפ"ד
% type: מבחן
% semester: א
% moed: מיוחד
% number: 7
% points: 14
% categories: improper-integrals, lhopital
% source: exams/מבחנים/תשפד/מועד ג תשפד.pdf

\begin{question}
חשבו את האיטגרל הלא-אמיתי הבא, או הראו התבדרות שלו:
\[ \int_2^\infty x\cdot e^{-x}\,dx \]
\end{question}

\begin{hint}
אינטגרציה בחלקים עם $u=x$, $v'=e^{-x}$, ואחר כך גבול כאשר הגבול העליון שואף לאינסוף.
\end{hint}

\begin{solution}
לפי ההגדרה, $\int_2^\infty xe^{-x}\,dx=\lim_{R\to\infty}\int_2^Rxe^{-x}\,dx$, אם הגבול קיים.

\textbf{פונקציה קדומה:} אינטגרציה בחלקים עם $u=x$, $v'=e^{-x}$ ($u'=1$, $v=-e^{-x}$):
\[ \int xe^{-x}\,dx=-xe^{-x}+\int e^{-x}\,dx=-xe^{-x}-e^{-x}+C=-(x+1)e^{-x}+C. \]
לכן
\[ \int_2^Rxe^{-x}\,dx=\Big[-(x+1)e^{-x}\Big]_2^R=3e^{-2}-\frac{R+1}{e^R}. \]
\textbf{מעבר לגבול:} $\lim_{R\to\infty}\frac{R+1}{e^R}$ הוא מקרה $\frac\infty\infty$, ולפי כלל לופיטל $\lim_{R\to\infty}\frac{1}{e^R}=0$. לכן
\[ \int_2^\infty xe^{-x}\,dx=\boxed{\frac{3}{e^2}}\approx0.406, \]
והאינטגרל מתכנס.
\end{solution}
